\section{Inleiding}%
\label{sec:inleiding}
Cybercriminaliteit die gericht is op het verkrijgen van geld of persoonsgegevens is geen recent verschijnsel. Reeds in de jaren negentig slaagden aanvallers erin om via phishing de authenticatiegegevens van America Online (AOL)-gebruikers te bemachtigen \autocite{Goel2018}. Sindsdien heeft phishing zich verder ontwikkeld en is het uitgegroeid tot een veelvoorkomende vorm van cybercriminaliteit. De blijvende relevantie van deze vorm van criminaliteit blijkt onder meer uit recente cijfers over de omvang van phishing in België.

In 2025 stuurden burgers bijna 10 miljoen verdachte berichten door naar \url{verdacht@safeonweb.be}\footnote{Centrum voor Cybersecurity België, \url{https://ccb.belgium.be/nl}, geraadpleegd op 8 oktober 2026.}. De omvang van phishing is bovendien niet uitsluitend zichtbaar in het aantal gemelde berichten, maar ook in de financiële schade die ermee gepaard gaat. Volgens \textcite{Febelfin2025} maakten fraudeurs in 2024 ongeveer 49 miljoen euro buit.

Hoewel phishing via verschillende kanalen verloopt, blijft misleiding via e-mail volgens \textcite{Veit2025} een van de grootste beveiligingsrisico's voor zowel bedrijven als eindgebruikers. Wanneer zo'n phishingmail de filters van de mailclient passeert, is het de verantwoordelijkheid van de gebruiker om in te schatten of de mail te vertrouwen is of niet. Deze beoordeling wordt bovendien bemoeilijkt door de opkomst van generatieve artificiële intelligentie, die aanvallers in staat stelt om overtuigende en gepersonaliseerde phishingberichten te creëren \autocite{SchmittFlechais2024}.

Hoewel sensibilisering gebruikers kan helpen om phishingmails te herkennen, volstaat ze niet om het risico op succesvolle aanvallen volledig weg te nemen. Onderzoek toont aan dat gebruikers ook na phishing\-training kwetsbaar blijven voor dergelijke aanvallen \autocite{Marshall2024}. 
Daarom is aanvullende ondersteuning wenselijk op het moment dat de gebruiker een e-mail moet beoordelen. 

Deze bachelorproef onderzoekt daarom de volgende centrale onderzoeksvraag: \emph{``Hoe kunnen relevante kenmerken uit e-mails worden geïdentificeerd en gecombineerd om phishingmails automatisch te detecteren en hun risiconiveau te bepalen?''}

Naast de scriptie levert het onderzoek een algoritme op dat een e-mail analyseert, er een risicoscore aan toekent en de onderliggende redenen op een begrijpelijke en toegankelijke manier aan de gebruiker toont. Om de centrale onderzoeksvraag te beantwoorden, worden de volgende deelvragen behandeld:

\begin{enumerate}
  \item \emph{Welke kenmerken zijn het meest relevant om phishing te detecteren?}
  \item \emph{Hoe kunnen deze kenmerken automatisch uit een e-mail worden gehaald?}
  \item \emph{Hoe kunnen deze kenmerken worden gecombineerd tot een risicoscore?}
\end{enumerate}

Het onderzoek wordt als geslaagd beschouwd wanneer het algoritme, geëvalueerd op een vooraf afgebakende dataset van phishingmails en legitieme e-mails, phishing betrouwbaar van legitieme e-mail onderscheidt, gemeten aan de hand van evaluatiecriteria die in de literatuurstudie worden onderbouwd.
